% Using InitechBase Chapter 5, entries SEEK .. ZAP, with the SET options
% (AN-04; SET options filed under SET, as in the model).

% SEEK -- os/samir/cmd/query.c q_seek (expression through ndx_seek on the
% master index; type must match the key; SET EXACT; #26).  Gates:
% test-ndx-seek, test-samir-query.
\begin{XbaseEntry}{SEEK}{command}
Finds the first record whose index key matches the value of an expression.
\Syntax
\SyntaxLine{\kw{SEEK} \xp{exp}}
\Usage
\cmd{SEEK} evaluates the expression and searches the master index of the
database in USE for it.  If the key of a record matches, the record pointer
moves to the first such record in index order and \fn{FOUND()} is true.  If
no key matches, the pointer moves to the end of the file: \fn{EOF()} is true
and \fn{FOUND()} false.

The expression must be of the type of the index key: a character expression
for an index on a character field, a number for an index on a numeric field.
For a character key, a key that begins with the value matches while SET EXACT
is OFF; with SET EXACT ON the key must equal the value, apart from trailing
blanks.

\cmd{SEEK} needs an open index; without one it displays
\xmsg{26}{Database is not indexed.}
\Tips
In a program, \cmd{SEEK} a memory variable that holds the answer to an
\cmd{ACCEPT}, then test \fn{FOUND()}.
\Examples
\DotExample{examples/c-seek}
The index \file{STAFFNO.NDX} is on the employee number:
\DotExample{examples/c-seek-2}
\SeeAlso{FIND, LOCATE, SET EXACT, SET ORDER, USE, FOUND()}
\end{XbaseEntry}

% SELECT -- os/samir/cmd/flow.c "SELECT" (a number 1..10 or an alias).
% Gate: test-interp-use.
\begin{XbaseEntry}{SELECT}{command}
Chooses the work area in which the following commands work.
\Syntax
\SyntaxLine{\kw{SELECT} \xp{n}/\xp{alias}}
\Usage
InitechBase has ten work areas, numbered 1 to 10, and each may hold one
database with its index files.  Work area 1 is selected when InitechBase
starts.  \cmd{SELECT} \xp{n} selects work area \xp{n}; \cmd{SELECT} \xp{alias}
selects the work area of the database with that alias.  The alias of a
database is its file name without the extension, unless \cmd{USE} gave it
another.

\cmd{USE} opens a database in the selected work area, and the commands that
read or change records work on the database of the selected work area.  Each
work area keeps its own record pointer, so moving in one database does not
move the others.
\Examples
\DotExample{examples/c-select}
\SeeAlso{USE, CLOSE, DBF()}
\end{XbaseEntry}

% SET CENTURY -- os/samir/cmd/set.c do_set_century (ctx->set_century; DTOC
% and display use four digits; CTOD accepts them).  Gate:
% test-interp-setfmt.
\begin{XbaseEntry}{CENTURY}{set}
Determines whether dates are displayed with the century.
\Syntax
\SyntaxLine{\kw{SET CENTURY} \kw{ON}/\kw{OFF}}
\Defaults
OFF.
\Usage
With SET CENTURY OFF, the year of a date has two digits: \texttt{12/25/91}.
With SET CENTURY ON, it has four: \texttt{12/25/1991}.  The setting applies to
dates displayed by \cmd{?}, \cmd{LIST} and \cmd{DISPLAY} and to the result of
\fn{DTOC()}.

A year typed with two digits is a year of the nineteen hundreds:
\fn{CTOD()} converts \texttt{12/25/91} to the 25th of December 1991 whatever the
setting.  With SET CENTURY ON, a four-digit year may also be given.
\Examples
\DotExample{examples/c-setcentury}
\SeeAlso{SET DATE, CTOD(), DTOC()}
\end{XbaseEntry}

% SET CONSOLE -- os/samir/cmd/set.c (ctx->set_console stored; no consumer in
% the shipped engine).
\begin{XbaseEntry}[the setting is recorded but the shipped version does not
  act on it: output is displayed with CONSOLE ON or OFF (set.c)]{CONSOLE}{set}
Turns the display of output on the screen on or off.
\Syntax
\SyntaxLine{\kw{SET CONSOLE} \kw{ON}/\kw{OFF}}
\Defaults
ON.
\Usage
InitechBase accepts SET CONSOLE ON and OFF and records the setting.  The
output of commands is displayed whatever the setting.
\SeeAlso{SET TALK, REPORT FORM}
\end{XbaseEntry}

% SET DATE -- os/samir/cmd/set.c do_set_date (eight formats; ctx->set_date_fmt
% read by DTOC/CTOD and display).  Gate: test-interp-setfmt.
\begin{XbaseEntry}{DATE}{set}
Chooses the form in which dates are displayed and typed.
\Syntax
\SyntaxLine{\kw{SET DATE} \opt{\kw{TO}} \kw{AMERICAN}/\allowbreak\kw{ANSI}/\allowbreak\kw{BRITISH}/\allowbreak\kw{FRENCH}/\allowbreak\kw{GERMAN}/\allowbreak\kw{ITALIAN}/\allowbreak\kw{JAPAN}/\allowbreak\kw{USA}}
\Defaults
AMERICAN.
\Usage
The setting determines the order of the month, day and year and the character
that separates them, both for dates that InitechBase displays and for the
character strings that \fn{CTOD()} converts:

\begin{tabular}{@{}ll@{}}
\kw{AMERICAN} & \texttt{mm/dd/yy}\\
\kw{ANSI} & \texttt{yy.mm.dd}\\
\kw{BRITISH} & \texttt{dd/mm/yy}\\
\kw{FRENCH} & \texttt{dd/mm/yy}\\
\kw{GERMAN} & \texttt{dd.mm.yy}\\
\kw{ITALIAN} & \texttt{dd-mm-yy}\\
\kw{JAPAN} & \texttt{yy/mm/dd}\\
\kw{USA} & \texttt{mm-dd-yy}\\
\end{tabular}

The year has four digits while SET CENTURY is ON.
\Examples
The first of March 1992 in each form:
\DotExample{examples/c-setdate}
\DotExample{examples/c-setdate-2}
\SeeAlso{SET CENTURY, CTOD(), DTOC(), DATE()}
\end{XbaseEntry}

% SET DECIMALS -- os/samir/cmd/set.c do_set_decimals (0..15, empty = 2;
% ctx->set_decimals read by display of computed values).  Gate:
% test-interp-decimals.
\begin{XbaseEntry}{DECIMALS}{set}
Sets the number of decimal places displayed for the results of calculations.
\Syntax
\SyntaxLine{\kw{SET DECIMALS TO} \opt{\xp{expN}}}
\Defaults
2.
\Usage
The setting determines how many decimal places \cmd{?}, \cmd{LIST} and
\cmd{DISPLAY} show for a value that InitechBase has calculated, such as the
result of a division.  The value itself keeps its full precision; only its
display is rounded.  Numeric fields keep the decimal places of their
structure, and whole numbers are displayed without decimal places.

The number of places may be from 0 to 15; a larger number gives 15.  SET
DECIMALS TO with no number restores 2.
\Examples
\DotExample{examples/c-setdecimals}
\SeeAlso{?, ROUND(), STR()}
\end{XbaseEntry}

% SET DELETED -- os/samir/cmd/set.c (wa_set_deleted: the walk of
% LIST/DISPLAY/LOCATE and navigation skips 0x2A records).  Gate:
% test-interp-set.
\begin{XbaseEntry}{DELETED}{set}
Determines whether records marked for deletion are processed.
\Syntax
\SyntaxLine{\kw{SET DELETED} \kw{ON}/\kw{OFF}}
\Defaults
OFF.
\Usage
With SET DELETED ON, the commands that move through the database and process
records -- \cmd{LIST}, \cmd{DISPLAY}, \cmd{LOCATE}, \cmd{CONTINUE},
\cmd{SKIP}, \cmd{GO TOP}, \cmd{GO BOTTOM} -- pass over the records marked for
deletion as if they were not there.  The records remain in the database until
\cmd{PACK}.  With SET DELETED OFF, marked records are processed like any
other, and are shown with an asterisk.
\Examples
\DotExample{examples/c-setdeleted}
\SeeAlso{DELETE, PACK, RECALL, DELETED()}
\end{XbaseEntry}

% SET EXACT -- os/samir/cmd/set.c do_set_exact (ctx->set_exact; eval.c
% C = C begins-with rule; ndx_seek).  Gates: test-interp-set,
% test-xbase-eval.
\begin{XbaseEntry}{EXACT}{set}
Determines how character strings are compared.
\Syntax
\SyntaxLine{\kw{SET EXACT} \kw{ON}/\kw{OFF}}
\Defaults
OFF.
\Usage
With SET EXACT OFF, the comparison \xp{expC1} \texttt{=} \xp{expC2} is true
when \xp{expC1} begins with \xp{expC2}: \texttt{'SMYKOWSKI' = 'SMY'} is true.
The string on the right decides how many characters are compared, so an empty
string on the right equals every string.  With SET EXACT ON, the two strings
must be the same, apart from trailing blanks.  \texttt{<>} and \texttt{\#}
follow the same rule, and so do \cmd{SEEK} and \cmd{FIND}.

The setting does not affect \texttt{<}, \texttt{>}, \texttt{<=}, \texttt{>=} or
the substring operator \texttt{\$}.  Capital and small letters are always
different.
\Examples
\DotExample{examples/c-setexact}
\SeeAlso{SEEK, FIND, AT()}
\end{XbaseEntry}

% SET FILTER -- os/samir/cmd/set.c do_set_filter -> wa_set_filter (the
% walk of LIST/DISPLAY/GO TOP/SKIP admits matching records only).  Gate:
% test-interp-set.
\begin{XbaseEntry}{FILTER}{set}
Makes InitechBase pass over the records that do not meet a condition.
\Syntax
\SyntaxLine{\kw{SET FILTER TO} \opt{\xp{condition}}}
\Usage
After SET FILTER TO \xp{condition}, the database of the selected work area
appears to hold only the records that meet the condition.  \cmd{LIST},
\cmd{DISPLAY}, \cmd{GO TOP}, \cmd{GO BOTTOM} and \cmd{SKIP} pass over the others.
\cmd{GO} with a record number still goes to that record.

The filter applies to the database in USE in the selected work area, until it
is removed with SET FILTER TO and no condition, or the database is closed.
SET FILTER needs a database in USE; without one it displays
\xmsg{17}{Cannot select requested database.}
\Tips
Move the record pointer with \cmd{GO TOP} after you set a filter, so that the
current record is one that meets it.
\Examples
\DotExample{examples/c-setfilter}
\SeeAlso{LOCATE, LIST, SET DELETED}
\end{XbaseEntry}

% SET HEADING -- os/samir/cmd/set.c (ctx->set_heading; query.c leaves out
% the heading line of LIST / DISPLAY).  Gates: test-interp-set,
% test-interp-list-heading-mutant.
\begin{XbaseEntry}{HEADING}{set}
Determines whether \cmd{LIST} and \cmd{DISPLAY} begin with a heading line.
\Syntax
\SyntaxLine{\kw{SET HEADING} \kw{ON}/\kw{OFF}}
\Defaults
ON.
\Usage
With SET HEADING ON, \cmd{LIST} and \cmd{DISPLAY} begin with a line that names
the columns.  With SET HEADING OFF, only the records are displayed.  Without
the headings, a column is as wide as its values only.
\Examples
\DotExample{examples/c-setheading}
\SeeAlso{LIST, DISPLAY}
\end{XbaseEntry}

% SET INDEX -- os/samir/cmd/set.c do_set_index (text recorded; the runtime
% open is deferred: "GATED: runtime index-open effect deferred").
\begin{XbaseEntry}[the index list is recorded but not opened by this version
  (set.c); open index files with USE ... INDEX]{INDEX}{set}
Names the index files to be used with the database in USE.
\Syntax
\SyntaxLine{\kw{SET INDEX TO} \opt{\xp{index file list}}}
\Usage
InitechBase accepts SET INDEX TO with or without a list of index files and
displays no message.  The index files open with the database are those named
in \cmd{USE} \ldots{} \kw{INDEX}.
\SeeAlso{USE, SET ORDER, SEEK}
\end{XbaseEntry}

% SET ORDER -- os/samir/cmd/set.c do_set_order -> wa_set_order (0 = natural
% order, n = the n-th index of USE ... INDEX).  Gate: test-interp-set.
\begin{XbaseEntry}{ORDER}{set}
Chooses which of the open index files is the master index.
\Syntax
\SyntaxLine{\kw{SET ORDER TO} \opt{\xp{n}}}
\Usage
When several index files are open with a database, the first named in
\cmd{USE} is the master index: it gives the order in which the records are
processed, and it is the index \cmd{SEEK} and \cmd{FIND} search.  SET ORDER
TO \xp{n} makes the \xp{n}th of the open index files the master index.  SET
ORDER TO 0, or SET ORDER TO alone, processes the records in record-number
order; the index files stay open.
\Tips
Move to the first record with \cmd{GO TOP} after SET ORDER.
\Examples
The open index files are on the name and on the department:
\DotExample{examples/c-setorder}
\SeeAlso{USE, SEEK, FIND, GO}
\end{XbaseEntry}

% SET RELATION -- os/samir/cmd/set.c do_set_relation (text recorded; the
% runtime effect is deferred: "GATED: runtime relation effect deferred").
\begin{XbaseEntry}[the relation is recorded but not maintained by this version
  (set.c)]{RELATION}{set}
Links the database in USE to a database in another work area.
\Syntax
\SyntaxLine{\kw{SET RELATION TO} \opt{\xp{key expression} \kw{INTO} \xp{alias}}}
\Usage
InitechBase accepts SET RELATION TO with or without a relation and displays no
message.  The record pointer of the other database does not move with the
selected one: select the other work area and use \cmd{SEEK} to find the related
record.
\SeeAlso{SELECT, SEEK, USE}
\end{XbaseEntry}

% SET SAFETY -- os/samir/cmd/set.c do_set_safety; os/samir/cmd/mutate.c
% m_confirm_file (ZAP asks "ZAP <file>? (Y/N)" while ON).  Gate:
% test-samir-zap-safety.
\begin{XbaseEntry}{SAFETY}{set}
Determines whether InitechBase asks for confirmation before it removes all
the records of a database.
\Syntax
\SyntaxLine{\kw{SET SAFETY} \kw{ON}/\kw{OFF}}
\Defaults
ON.
\Usage
With SET SAFETY ON, \cmd{ZAP} names the database and asks \texttt{(Y/N)}
before it removes the records.  Only the answer \type{Y}, in either case,
lets it go on; any other answer leaves the database as it is.  With SET
SAFETY OFF, \cmd{ZAP} removes the records without asking.
\Examples
See \cmd{ZAP}.
\SeeAlso{ZAP, PACK}
\end{XbaseEntry}

% SET STATUS -- os/samir/cmd/set.c (ctx->set_status) ->
% samir_status_refresh -> pal_milton.c milton_status / core/screen.c
% samir_screen_status (row 23: "Command Line", drive, database, Rec: n/m;
% the dot prompt scrolls in rows 1-22).  Gates: test-samir-status,
% test-samir-screen-emu.
\begin{XbaseEntry}{STATUS}{set}
Displays or removes the status bar at the bottom of the screen.
\Syntax
\SyntaxLine{\kw{SET STATUS} \kw{ON}/\kw{OFF}}
\Defaults
OFF.
\Usage
With SET STATUS ON, InitechBase clears the screen and keeps the two bottom
lines for itself.  The status bar, in reverse video, shows \texttt{Command
Line}, the current drive, the name of the database in USE in the selected
work area, and the current record number with the number of records
(\texttt{Rec: 3/13}).  The line below it is a guidance line.  Commands and their
output scroll in the lines above.  The status bar is brought up to date after
every command.

SET STATUS OFF removes the status bar; output then uses the whole screen.
\cmd{CLOSE} and \cmd{QUIT} also remove it.  The status bar appears only on the
screen; it is never part of output that is redirected to a file.
\SeeAlso{RECNO(), RECCOUNT(), DBF()}
\end{XbaseEntry}

% SET TALK -- os/samir/cmd/set.c do_set_talk (ss->talk stored; no consumer
% in the shipped engine).
\begin{XbaseEntry}[the setting is recorded but the shipped version does not
  act on it (set.c)]{TALK}{set}
Determines whether InitechBase reports the results of commands on the screen.
\Syntax
\SyntaxLine{\kw{SET TALK} \kw{ON}/\kw{OFF}}
\Defaults
ON.
\Usage
InitechBase accepts SET TALK ON and OFF and records the setting.  The commands
of this version display no result messages, and their output is the same with
either setting.  Programs commonly begin with SET TALK OFF.
\SeeAlso{SET CONSOLE}
\end{XbaseEntry}

% SKIP -- os/samir/cmd/query.c q_skip (n records in master-index order;
% stops at EOF/BOF without a message).  Gate: test-interp-nav.
\begin{XbaseEntry}{SKIP}{command}
Moves the record pointer forward or back.
\Syntax
\SyntaxLine{\kw{SKIP} \opt{\xp{expN}}}
\Defaults
1.
\Usage
\cmd{SKIP} moves the record pointer \xp{expN} records forward, or back if
\xp{expN} is negative.  With an index open, the records are counted in index
order; while a filter is set or SET DELETED is ON, the records passed over do
not count.

Moving forward past the last record puts the record pointer at the end of the
file: \fn{RECNO()} is one more than the number of records and \fn{EOF()} is
true.  Moving back past the first record leaves it on the first record, with
\fn{BOF()} true.  InitechBase displays no message in either case.
\Examples
\DotExample{examples/c-skip}
\SeeAlso{GO, EOF(), BOF(), RECNO()}
\end{XbaseEntry}

% STORE and = -- os/samir/cmd/flow.c "STORE" and "<name> = <expr>"
% (xb_interp_store: create or update, type from the value).  Gate:
% test-interp-flow.
\begin{XbaseEntry}{STORE}{command}
Stores a value in one or more memory variables.
\Syntax
\SyntaxLine{\kw{STORE} \xp{exp} \kw{TO} \xp{memvar list}}
\par or\par
\SyntaxLine{\xp{memvar} \kw{=} \xp{exp}}
\Usage
\cmd{STORE} evaluates the expression and stores its value in each memory
variable of the list.  A memory variable that does not exist is created; one
that exists receives the new value, whatever its type was.  The memory
variable takes the type of the value: character, numeric, date or logical.

The second form stores the value in one memory variable.  It does the same as
\cmd{STORE}.

A memory variable name has up to 10 letters, digits and underscores, and
begins with a letter.  Do not give a memory variable the name of a field of the
database in USE: in an expression the field of the current record is found
first, and the memory variable cannot be reached.  The programs of this manual
begin the names of their memory variables with the letter M.
\Examples
\DotExample{examples/c-store}
\SeeAlso{RELEASE, PUBLIC, PRIVATE, ACCEPT, INPUT, REPLACE}
\end{XbaseEntry}

% USE -- os/samir/samir_main.c repl_do_use (default .DBF / .NDX
% extensions; ALIAS; INDEX up to 7; first index master; first record in
% index order; #1, #3, #15, #28, #29).  Gates: test-samir-repl,
% test-interp-use, test-samir-use.
\begin{XbaseEntry}{USE}{command}
Opens a database, with its index files, in the selected work area, or closes
the database there.
\Syntax
\SyntaxLine{\kw{USE} \opt{\xp{file} \opt{\kw{ALIAS} \xp{alias}}
  \opt{\kw{INDEX} \xp{index file list}}}}
\Usage
\cmd{USE} \xp{file} closes the database open in the selected work area, if
there is one, and opens the database \xp{file}.  The extension \file{.DBF} is
added when none is given.  If the database has memo fields, its memo file, of
the same name with the extension \file{.DBT}, is opened with it.  The first
record becomes the current record.

\kw{INDEX} opens up to seven index files with the database.  The extension
\file{.NDX} is added when none is given.  The first index file of the list is
the master index: it gives the order of the records, and the first record in
that order becomes the current record.  The other index files are kept up to
date when records are added or changed.

\kw{ALIAS} gives the database another name for \cmd{SELECT} and \fn{DBF()}.
Without it, the alias is the file name without the extension.

\cmd{USE} with no file name closes the database of the selected work area.

\cmd{USE} reports a database that cannot be opened: \xmsg{1}{File does not
exist.} when there is no such file, \xmsg{3}{File is already open.} when it is
open in another work area, and \xmsg{15}{Not a dBASE database.} when the file
is not a database.

\cmd{USE} is a dot prompt command.  A program works on the databases that are
open when it is run.
\Examples
\DotExample{examples/c-use}
\DotExample{examples/c-use-2}
\DotExample{examples/c-use-3}
\SeeAlso{CLOSE, SELECT, SET ORDER, DBF()}
\end{XbaseEntry}

% WAIT -- os/samir/cmd/proc.c p_wait (prompt, one key via conin_char, a
% one-character string; Enter alone -> "").  Gate: test-interp-proc.
\begin{XbaseEntry}{WAIT}{command}
Waits for one key to be pressed, and can store it in a memory variable.
\Syntax
\SyntaxLine{\kw{WAIT} \opt{\xp{expC}} \opt{\kw{TO} \xp{memvar}}}
\Usage
\cmd{WAIT} displays the prompt \xp{expC}, if one is given, and waits until a key
is pressed.  The key is not followed by \key{Enter}.  With \kw{TO}, the key is
stored in \xp{memvar} as a character string of one character; if the key was
\key{Enter}, the memory variable is an empty string.  Without a prompt,
\cmd{WAIT} displays nothing while it waits.
\Tips
A program uses \cmd{WAIT} to hold its output on the screen until the operator
has read it, and to take a single-letter answer to a question.
\Examples
\DotExample{examples/c-wait}
\SeeAlso{ACCEPT, INPUT}
\end{XbaseEntry}

% ZAP -- os/samir/cmd/mutate.c m_zap (m_confirm_file while SAFETY ON; the
% table emptied in place, its indexes emptied).  Gates:
% test-samir-zap-safety, test-interp-replace.
\begin{XbaseEntry}{ZAP}{command}
Removes all the records of the database in USE.
\Syntax
\SyntaxLine{\kw{ZAP}}
\Usage
\cmd{ZAP} removes every record from the database of the selected work area and
leaves an empty database with the same structure.  The records cannot be
recovered.  Use \cmd{ZAP} on a database opened without index files.

While SET SAFETY is ON, \cmd{ZAP} first asks \texttt{ZAP} \xp{file}\texttt{? (Y/N)};
the records are removed only if the answer is \type{Y}.
\begin{AteWarning}
\cmd{ZAP} cannot be undone.  Keep a copy of a database before you remove its
records.
\end{AteWarning}
\Examples
\DotExample{examples/c-zap}
With SET SAFETY OFF the records are removed at once:
\DotExample{examples/c-zap-2}
\SeeAlso{DELETE, PACK, SET SAFETY}
\end{XbaseEntry}
